% Preprint-style rendering of the RNAAS note (article class), for viewing
% where aastex631.cls is unavailable. The submission source is
% spetc_rnaas_aastex.tex; the scientific text is identical.
\documentclass[11pt, a4paper]{article}
\usepackage{fontspec}
\setmainfont{TeX Gyre Termes}
\setsansfont{TeX Gyre Heros}
\setmonofont{TeX Gyre Cursor}
\usepackage[final]{microtype}
\usepackage{amsmath, amssymb, mathtools}
\usepackage{graphicx}
\usepackage[margin=2.6cm]{geometry}
\usepackage{natbib}
\bibliographystyle{plainnat}
\usepackage[colorlinks=true, linkcolor=blue, citecolor=blue,
            urlcolor=blue, filecolor=blue]{hyperref}

\title{SPETC: An Open-Source Spectro-Photometric Exposure Time Calculator
for Small Telescopes and Amateur Instrumentation\\[4pt]
{\normalsize Research Note of the AAS --- preprint rendering}}
\author{Mauro Barbieri\\\small\texttt{mauro.barbieri@pm.me}}
\date{}

\begin{document}
\maketitle

\noindent\textbf{Keywords:} Astronomy software (1855) --- Observational
astronomy (1145) --- CCD photometry (208) --- Spectroscopy (1558) ---
Amateur astronomy (35)

\begin{abstract}
SPETC is an open-source, offline exposure time calculator for ground-based
broad-band photometry and slit or slitless spectroscopy, aimed at the
instrumentation actually used on small telescopes: CMOS cameras with
gain-dependent read noise, one-shot-colour sensors, classical slit
spectrographs, and Star Analyser transmission gratings. It couples
CALSPEC-calibrated template spectra and SVO-convention synthetic photometry
to a position- and time-dependent sky model (airglow, zodiacal light,
starlight, Krisciunas--Schaefer moonlight, twilight), a complete noise
budget including scintillation and sky-subtraction noise, and per-night
planning tracks with refraction-corrected airmass and parallactic angle.
Every relation is documented, unit-tested, and inverted in closed form for
target-S/N planning.
\end{abstract}

\section{Motivation}

Exposure time calculators are standard infrastructure at professional
observatories, but each is bound to its facility, and the generic
alternatives model idealized instruments. The communities producing a
growing share of time-domain follow-up --- variable-star and exoplanet
photometrists, Be-star and nova spectroscopists on sub-metre telescopes ---
plan observations with rules of thumb, although their regime is the one
where honest error budgets matter most: bright stars are
scintillation-limited, urban skies are background-limited, and slitless
gratings are sky-limited in a way slit intuition does not predict. SPETC
(Spectro-Photometric Exposure Time Calculator) addresses exactly this
regime. It descends from a 2007 Fortran chain, rewritten in Python on
\texttt{astropy} \citep{astropy2022}, and runs fully offline as a desktop
application, a headless batch tool, and an importable library.

\section{Physics}

Templates are absolutely calibrated CALSPEC/BPGS spectrophotometry
\citep{bohlin2014}, rescaled to the user's magnitude in any reference
filter through SVO-convention synthetic photometry \citep{rodrigo2020}
with photon or energy weighting, and optionally transformed by a radial
velocity and CCM89 reddening \citep{cardelli1989}. The sky model is
position- and time-dependent: van Rhijn airglow attenuated along the
slant path, zodiacal light as a function of ecliptic latitude and solar
elongation, integrated starlight, \citet{krisciunas1991} moonlight with
the lunar distance factor, and a twilight/daylight bridge on the
geometric solar altitude. Airmass follows \citet{pickering2002} on
refracted apparent altitudes; atmospheric dispersion follows
\citet{filippenko1982} and feeds worst-case slit losses whenever the slit
is not at the parallactic angle, which the planning track reports per
time sample. Point-spread functions are Gaussian or Moffat with analytic
encircled-energy and slit couplings (the Moffat marginal is a Student-$t$
distribution, so no profile is ever numerically integrated), with
optional Kolmogorov seeing scaling in wavelength and airmass and a
guiding-blur term. The noise budget is the CCD equation extended with
scintillation \citep{young1967,osborn2015}, ADC quantization, dark, an
extra-background term, and the \citet{merline1995} sky-subtraction
factor; because the scintillation variance is linear in time, the
target-S/N inversion and a stack planner (frames $\times$ sub-exposure,
read-noise penalty, sky-limited frame length) remain closed-form even
for scintillation-limited stars. Spectroscopic counts are exact per-resel
integrals of the template --- narrow-line fluxes are invariant under the
chosen resolving power --- and each output row carries the
\citet{cayrel1988} equivalent-width uncertainty. Slitless mode implements
Star Analyser physics: dispersion from groove density and grating-sensor
distance, seeing-limited resolution, and the full-bandpass sky background
that every pixel of a slitless spectrum receives, the dominant and often
overlooked term. Figure~\ref{fig:budget} shows the resulting behaviour
for a representative 358-mm configuration.

\section{Implementation and availability}

SPETC provides a Tkinter desktop interface with whole-night planning
plots, a JSON-driven batch mode producing CSV and HTML summaries, tools
to import CALSPEC/STScI FITS spectra as templates and to convert measured
filter transmissions into fully annotated SVO-format profiles, CMOS
gain-table support, SQM/Bortle sky input, differential-photometry
precision in mmag, and a terrain-horizon generator based on the
Copernicus GLO-30 elevation model. All physics is regression-tested
(70+ tests), documented in a derivation-level manual, and was subjected
to a full internal physics audit whose findings --- including the
slitless sky term and the resolution-invariance requirement --- are
implemented and tested. The code, documentation and audit are available
under an open-source license in the project repository.

\begin{figure}[ht!]
\centering
\includegraphics[width=0.7\textwidth]{figures/fig_06_snr_budget.png}
\caption{Broad-band signal-to-noise in 60 s versus target magnitude for a
358-mm telescope, computed by the released code. The full budget (solid)
saturates at the scintillation ceiling near $\mathrm{S/N} \simeq 1.2
\times 10^{3}$; a Poisson-only model (dashed) overestimates bright-star
performance by more than an order of magnitude --- the central planning
error SPETC is designed to prevent.}
\label{fig:budget}
\end{figure}

\bibliography{references}

\end{document}
